\section{The Kitaev chain}\label{sec:kitaev_chain}
%% ===================================================================

The Kitaev chain, or Majorana chain~\cite{Kitaev2001Majorana}, is a chain of
spinless fermions. In the limit where the Majorana modes pair up only across
neighbouring sites, its ground state is the non-trivial fixed point of the
topological phase, and it is described by a graded (fermionic) matrix product
state with bond dimension two and a twist at the boundary
\cite[Appendix~A, ``The Kitaev chain'']{Cirac2021Matrix},
\cite{Bultinck2017FermionicMPS}. On a ring of $N$ sites the state is
\begin{align}
    \ket\psi
     & =
    \sum_{i_1,\ldots,i_N}\tr\bigl(YA^{i_1}A^{i_2}\cdots A^{i_N}\bigr)
    \ket{i_1}\otimes_g\ket{i_2}\otimes_g\cdots\otimes_g\ket{i_N},
    \label{eq:kitaev_chain_state}
\end{align}
where $\otimes_g$ is the graded tensor product; the state has odd fermion
parity~\cite[Section~II.B.4]{Cirac2021Matrix}.

\paragraph{Scope.} Graded tensor networks are not formalized. The results below
concern the ordinary contraction of the matrices, which gives the coefficients
of~\eqref{eq:kitaev_chain_state} in the ordered graded product basis; the
graded contraction rules and the graded notion of injectivity are outside
their scope~\cite{gap:rmp_kitaev_chain_bosonic_scope}.

\begin{definition}[Kitaev-chain tensor and twist]\label{def:kitaev_tensor}
    \lean{MPSTensor.kitaevTensor, MPSTensor.kitaevTwist}
    \leanok
    \uses{def:mps_tensor}
    Set $d=D=2$. The Kitaev-chain tensor is
    \begin{align}
        A^0 & = \begin{pmatrix}1&0\\0&1\end{pmatrix},
        \qquad
        A^1 = \begin{pmatrix}0&1\\-1&0\end{pmatrix},
        \notag
    \end{align}
    with the twist $Y=\begin{pmatrix}0&1\\-1&0\end{pmatrix}$ at the boundary
    \cite[Appendix~A, ``The Kitaev chain'']{Cirac2021Matrix},
    \cite{Bultinck2017FermionicMPS}.
\end{definition}

\begin{definition}[Twisted periodic amplitude]\label{def:kitaev_twisted_coeff}
    \lean{MPSTensor.kitaevTwistedCoeff}
    \leanok
    \uses{def:kitaev_tensor, def:eval_word}
    For a word $w=(i_1,\ldots,i_N)$ the twisted periodic amplitude is
    $\tr(YA^{i_1}\cdots A^{i_N})$, the coefficient of $w$
    in~\eqref{eq:kitaev_chain_state}.
\end{definition}

\begin{lemma}[The twist squares to minus the identity]\label{lem:kitaev_twist_sq}
    \lean{MPSTensor.kitaevTensor_one, MPSTensor.kitaevTwist_mul_self,
        MPSTensor.kitaevTensor_one_mul_self, MPSTensor.kitaevTwist_mul_kitaevTensor}
    \leanok
    \uses{def:kitaev_tensor}
    $A^1=Y$, $(A^1)^2=Y^2=-\mathbb 1$, and $Y$ commutes with $A^0$ and $A^1$
    \cite{Bultinck2017FermionicMPS}.
\end{lemma}

\begin{proof}
    \leanok
    Direct multiplication of $2\times2$ matrices.
\end{proof}

\begin{lemma}[Products of the Kitaev-chain matrices]\label{lem:kitaev_eval_word}
    \lean{MPSTensor.kitaevTensor_evalWord}
    \leanok
    \uses{def:kitaev_tensor, def:eval_word, lem:kitaev_twist_sq}
    If the word $w$ contains the letter $1$ exactly $k$ times, then
    $A^{i_1}\cdots A^{i_N}=Y^k$.
\end{lemma}

\begin{proof}
    \leanok
    Induction on the word, using $A^0=\mathbb 1$ and $A^1=Y$.
\end{proof}

\begin{theorem}[Periodic amplitudes]\label{thm:kitaev_amplitudes}
    \lean{MPSTensor.kitaevTensor_coeff, MPSTensor.kitaevTwistedCoeff_eq}
    \leanok
    \uses{def:mpv, def:kitaev_twisted_coeff, lem:kitaev_eval_word}
    Let the word $w$ contain the letter $1$ exactly $k$ times. Then
    \begin{align}
        \tr(A^{i_1}\cdots A^{i_N})
         & =
        \begin{cases}
            2(-1)^{k/2}, & k \text{ even}, \\
            0,           & k \text{ odd},
        \end{cases}
        \qquad
        \tr(YA^{i_1}\cdots A^{i_N})
        =
        \begin{cases}
            0,               & k \text{ even}, \\
            2(-1)^{(k+1)/2}, & k \text{ odd}.
        \end{cases}
        \notag
    \end{align}
    Neither amplitude depends on the length $N$ except through $k$.
\end{theorem}

\begin{proof}
    \leanok
    By the preceding lemma the products are $Y^k$ and $Y^{k+1}$. Since
    $Y^2=-\mathbb 1$, an even power $Y^{2m}$ is $(-1)^m\mathbb 1$, of trace
    $2(-1)^m$, and an odd power $Y^{2m+1}$ is $(-1)^mY$, of trace zero.
\end{proof}

\begin{corollary}[Complementary parity supports]\label{cor:kitaev_parity_supports}
    \lean{MPSTensor.kitaevTwistedCoeff_ne_zero_iff,
        MPSTensor.kitaevTensor_coeff_ne_zero_iff,
        MPSTensor.kitaevTensor_coeff_mul_kitaevTwistedCoeff}
    \leanok
    \uses{thm:kitaev_amplitudes}
    The twisted amplitude of a word is nonzero exactly when the word contains
    an odd number of letters $1$, in agreement with the odd parity of the
    state~\eqref{eq:kitaev_chain_state}; the untwisted amplitude is nonzero
    exactly when that number is even. In particular the two amplitudes of a
    word never both differ from zero.
\end{corollary}

\begin{proof}
    \leanok
    Read off from the closed forms.
\end{proof}

\begin{lemma}[Open-boundary closures]\label{lem:kitaev_open_boundary}
    \lean{MPSTensor.kitaevTensor_evalWord_entries}
    \leanok
    \uses{lem:kitaev_eval_word}
    For every word, the product $M=A^{i_1}\cdots A^{i_N}$ satisfies
    $M_{00}=M_{11}$ and $M_{01}=-M_{10}$: of the four ways to close the
    virtual indices at the two ends of an open chain, the two diagonal
    closures give the same amplitudes and the two off-diagonal closures give
    amplitudes differing by a sign, in the ordinary contraction
    \cite{Bultinck2017FermionicMPS}.
\end{lemma}

\begin{proof}
    \leanok
    Both $\mathbb 1$ and $Y$ have this form, and so does every power of $Y$.
\end{proof}

\begin{theorem}[The ordinary contraction is reducible and not normal]%
    \label{thm:kitaev_not_normal}
    \lean{MPSTensor.kitaevTensor_mulVec_common_eigenvector,
        MPSTensor.kitaevTensor_not_isInjective, MPSTensor.kitaevTensor_not_isNormal}
    \leanok
    \uses{def:kitaev_tensor, def:injective, def:normal, lem:kitaev_open_boundary}
    Read as an ordinary matrix product state, the Kitaev-chain tensor is
    reducible: the vector $(1,i)$ is a common eigenvector of $A^0$ and $A^1$,
    with eigenvalues $1$ and $i$. The tensor is not injective, and it is not
    normal: no blocking length makes the products span the full matrix
    algebra. This is the bosonic case that
    \cite[Section~II.B.4]{Cirac2021Matrix} contrasts with the injectivity of
    the graded description.
\end{theorem}

\begin{proof}
    \leanok
    The linear functional $M\mapsto M_{00}-M_{11}$ vanishes on every product
    by the preceding lemma, and it does not vanish on the matrix unit
    $\ketbra{0}{0}$, so the products of any fixed length span a proper
    subspace.
\end{proof}

The Kitaev-chain tensor is related to the even-parity tensor
$E^0=\tfrac1{\sqrt2}\mathbb 1$, $E^1=\tfrac1{\sqrt2}\sigma^x$, whose periodic
contraction is the equal superposition of even-parity configurations.

\begin{lemma}[Relation to the even-parity tensor]\label{lem:kitaev_even_parity_bridge}
    \lean{MPSTensor.kitaevGauge, MPSTensor.kitaevGauge_mul_inv,
        MPSTensor.kitaevTensor_mul_kitaevGauge}
    \leanok
    \uses{def:kitaev_tensor}
    Let $X=\operatorname{diag}(1,i)$, which is invertible with inverse
    $\operatorname{diag}(1,-i)$, and let $\varphi_0=1$, $\varphi_1=i$. Then
    $A^jX=\sqrt2\,\varphi_j\,XE^j$ for $j=0,1$: up to the on-site phase
    $\operatorname{diag}(1,i)$ on the physical index and the factor $\sqrt2$,
    the Kitaev-chain tensor is gauge equivalent to the even-parity tensor.
\end{lemma}

\begin{proof}
    \leanok
    $X\sigma^xX^{-1}=\sigma^y$ and $i\sigma^y=Y$.
\end{proof}
